



\section{Automatic Evaluation}




\vspace{-3pt}
Our benchmark measures the following three dimensions of system responses:
\vspace{-3pt}
\begin{itemize}[leftmargin=1em,itemsep=0pt]
    \item \tf{Fluency}: whether the model's generated text is fluent and coherent.
    \item \tf{Correctness}: whether the answer is accurate and covers all aspects of interest.
    \item \tf{Citation quality}: whether the answer is well supported by the cited passages and no irrelevant passages are cited.
\end{itemize}
\vspace{-3pt}
In the following, we present automatic metrics for each dimension 
and discuss why the combination of the three metrics provides a robust evaluation. %



\subsection{Fluency} 
We use MAUVE \cite{pillutla2021mauve} to evaluate the fluency of the output (\S\ref{app_sec:imp_details}).
We deploy MAUVE for ASQA and ELI5
and omit it for QAMPARI,
as QAMPARI 
only requires a list of short answers as the response
and LLMs consistently adhere to the format  in our experiments.
As MAUVE is sensitive to output length and text style, and most LLMs are capable of producing fluent text,
we mainly employ it as a sanity check as long as the MAUVE scores are high enough.



\subsection{Correctness}
\label{sec:correctness}





Our objective is to measure the informativeness and utility of the generation  to the question. 
\citet{liu2023evaluating} propose to directly evaluate \textit{perceived utility}
by humans, a process difficult to automate. 
Therefore, we use correctness---whether the response is accurate compared to a ground truth answer---as a proxy.
Evaluating the correctness of long-form generation is a challenging task~\citep{krishna-etal-2021-hurdles},
and we describe our strategy for each dataset below.
Figure~\ref{fig:correctness} illustrates the metrics and we include additional implementation details in \S\ref{app_sec:imp_details}.


For \tf{ASQA}, we follow \citet{stelmakh2022asqa} and 
calculate the recall of correct short answers
by checking whether the short answers (provided by the dataset) are exact substrings of the generation (\textit{exact match recall}; EM recall).


For \tf{QAMPARI}, we follow \citet{rubin2022qampari} and calculate the \textit{precision} and \textit{recall} of 
the model prediction, 
by checking the exact match to the gold answer list.
We add one additional adjustment:
considering that users often  want to know only a few example answers of 
the question, %
our evaluation considers recall to be 100\% if the prediction includes at least 5 correct answers (\textit{recall-5}).

Unlike ASQA and QAMPARI, the \tf{ELI5} dataset does not provide short entity answers.
\citet{fan-etal-2019-eli5} use ROUGE for evaluation, which does not reflect the correctness  well (\citealp{krishna-etal-2021-hurdles}; \S\ref{app_sec:claim_eli5}).
Inspired by works in summarization evaluation~\cite{zhang-bansal-2021-finding,kamoi2023wice,wang-etal-2020-asking}, 
we use InstructGPT~(\ttt{text-davinci-003}; \citealp{ouyang2022training}) to generate three ``sub-claims''. 
Then we use TRUE\footnote{\url{https://huggingface.co/google/t5_xxl_true_nli_mixture}. Details in \S\ref{app_sec:imp_details}.}~\cite{honovich-etal-2022-true-evaluating}, a T5-11B~\citep{raffel2020exploring} model fine-tuned on a collection of natural language inference (NLI)  datasets, to check whether the model output entails the sub-claims (\textit{claim recall}). TRUE targets factual correctness and has been used by previous works in similar context~\citep{bohnet2022attributed,gao2022rarr}.
We demonstrate that claim recall provides a more accurate measure of correctness than existing metrics (more details in \S\ref{app_sec:claim_eli5}).



\begin{figure}[t]
    \centering
    \includegraphics[width=0.94\linewidth]{figs/correctness.pdf}
    \caption{Evaluation of correctness (details in \S\ref{sec:correctness}).}
    \vspace{-13pt}
    \label{fig:correctness}
\end{figure}


\subsection{Citation Quality}
\label{sec:citation}





We evaluate citation qualities using two metrics:  
(1) \textit{citation recall}, which determines if the output is entirely supported by cited passages, and
(2) \textit{citation precision}, which identifies any irrelevant citations. %
Although we prioritize citation recall as it entails a well-supported and truthful answer, 
enhancing precision is crucial for better user satisfaction,
reducing the need for human review of extraneous passages.
Figure~\ref{fig:citation} provides an illustrated example. %

\begin{figure}[t]
    \centering
    \includegraphics[width=0.98\linewidth]{figs/citation.pdf}
    \caption{Evaluation of citation quality (details in \S\ref{sec:citation}). We use 
   an NLI model to verify whether a statement is supported by its citations. %
    }
    \vspace{-15pt}
    \label{fig:citation}
\end{figure}


We use the NLI model TRUE~\citep{honovich-etal-2022-true-evaluating} again
to automatically examine whether the cited passages entail the model generation.
We conduct human evaluation (\S\ref{sec:human}) to demonstrate strong human correlation of our metric.





\paragraph{Citation recall.} 
We calculate the citation recall of \textit{each statement} (0 or 1) and average
over all {statements} in the model response. %
For each statement $s_i$, %
its citation recall is $1$ if and only if there is at least one citation ($\mathcal{C}_i\neq \emptyset$) and 
$\phi(\text{concat}(\mathcal{C}_i), s_i)=1$,
where 
$\phi(\text{premise},\text{hypothesis})$
is the NLI model that outputs $1$ if the premise  entails the hypothesis, and $0$ otherwise; 
$\text{concat}(\mathcal{C}_i)$ concatenates all passages in $\mathcal{C}_i$ together (details in \S\ref{app_sec:imp_details}).
The NLI evaluation is in accordance with the \textit{attributable to identified sources} (AIS) framework~\cite{rashkin2021measuring}: 
$\phi(\text{concat}(\mathcal{C}_i), s_i)=1$ implies that 
$s_i$ is true based solely on $\text{concat}(\mathcal{C}_i)$.



 



\paragraph{Citation precision.} 
Our citation precision evaluation detects citations that are irrelevant,
but it does not require citing a minimal set.
We follow this design  because human writing often cites  redundant sources to enhance  credibility;
human readers may also appreciate multiple citations, especially when it pertains to critical claims such as medical advice.


We calculate the citation precision for \textit{each citation} (0 or 1) and
average over all {citations} in the response. %
We first define if a citation is ``irrelevant''.
Intuitively, a citation $c_{i,j}$ is ``irrelevant'' if (a) $c_{i,j}$ itself cannot support $s_i$ and (b) removing $c_{i,j}$ does not affect the rest of the citations to support $s_i$. Formally, $c_{i,j}$ is ``irrelevant'' if and only if
\begin{enumerate}[itemsep=0pt,leftmargin=2em,label=(\alph*)]
        \item  $\phi(c_{i,j},s_i)=0$, AND
        \item  $\phi(\text{concat}(\mathcal{C}_i \setminus \{c_{i,j}\}), s_i)=1$.  
\end{enumerate}
$c_{i,j}$ has a precision  of 1 if $s_i$ has recall$=$$1$ and $c_{i,j}$ is not irrelevant.
For example (Figure~\ref{fig:citation}), when $s_3$ cites three references \ttt{[2][4][5]} and recall$=$$1$,
\ttt{[2]} is ``irrelevant'' if $\phi(\texttt{[2]}, s_3)=0$ and $\phi(\texttt{[4][5]}, s_3)=1$.
For condition (b) to work, we set recall$=$$1$ as a prerequisite for precision$=1$. Note that this algorithm overlooks the scenario when one citation partially supports the statement. 
We discuss the details in \S\ref{app_sec:citation_recall}.









\subsection{\citeeval{} is Robust to Shortcut Cases}
We showcase how the \citeeval{} evaluation is robust to two possible shortcuts in \S\ref{app_sec:short_cut}:
(1) using the top-1 retrieved passage as the response and citing itself, 
and (2) using the first two sentences of the top-1 passage.
Both cases have almost-perfect citation scores, but (1) has low fluency due to its unnaturally long length compared to human answers, and (2) has low correctness due to low coverage.




















 